\documentclass[10pt]{article}
\usepackage[cmex10]{amsmath}
\usepackage{amssymb,amsthm}
\DeclareFontShape{OMX}{cmex}{m}{n}{<->cmex10}{}
\DeclareMathSizes{10}{10}{7}{7}
\DeclareMathSizes{9}{10}{7}{7}
\usepackage[margin=1in]{geometry}
\usepackage[hidelinks]{hyperref}
\newtheorem{theorem}{Theorem}
\title{Exact certificates for list edge weighting on small graphs}
\author{}
\date{November 2026}
\begin{document}
\maketitle

\begin{abstract}
We verify total $(1,3)$-weight choosability for all finite simple graphs whose
components have at most nine vertices and none is isomorphic to $K_2$. Every connected
case has an explicit coefficient certificate over a prime field. We also give
an exact fixed-graph search for lists that prevent a proper weighting.
\end{abstract}

\begin{theorem}
Let $G=(V,E)$ be a finite simple graph whose connected components have at
most nine vertices and none is isomorphic to $K_2$.
Assign a real number $a_v$ to each vertex and a set $L_e$ of three
distinct real weights to each edge. There is a choice $x_e\in L_e$ such that
$a_u+\sum_{e\ni u}x_e\ne a_v+\sum_{e\ni v}x_e$ for every $uv\in E$.
\end{theorem}

\begin{proof}
The components can be treated separately, and isolated vertices impose no
constraints. For a connected component with $m$ edges, consistently order
the endpoints of each edge and put $S_v=\sum_{e\ni v}x_e$. Use the polynomial
\[
P_a(x)=\prod_{uv\in E}(S_u-S_v+a_u-a_v).
\]
Its homogeneous part of degree $m$ is $P_0$. By the Combinatorial
Nullstellensatz \cite{alon}, any nonzero coefficient $[x^\alpha]P_0$ with
$\sum_e\alpha_e=m$ and $\alpha_e\leq 2$ certifies a choice from every such
assignment of edge lists, for all offsets $a_v$. This is the algebraic
criterion used in weight choosability \cite{bgn,zhu}.

For each catalogue graph, the supplementary data contain such an $\alpha$
and the coefficient residue modulo $p=1\,000\,000\,007$. For each tuple
$b$ of integers with $0\leq b_e\leq \alpha_e$, compute $P_0(b)$ by taking
the vertex sums of the edge weights $b_e$. The exact coefficient is
\[
[x^\alpha]P_0=
\frac{1}{\prod_e\alpha_e!}
\sum_{0\leq b\leq\alpha}(-1)^{m-\sum_e b_e}
\prod_e\binom{\alpha_e}{b_e}\,P_0(b).
\]
Indeed, the numerator is the mixed forward difference of orders $\alpha$ at
zero. It annihilates any monomial with exponent below $\alpha_e$ in some
coordinate. The total degree is $m$, so the only monomial left is
$x^\alpha$, which contributes $\prod_e \alpha_e!$ times its coefficient.
The denominator is a power of two, so the formula also holds modulo $p$.
Every archived residue is nonzero, which proves the corresponding integer
coefficient is nonzero.

We checked all connected graphs in the atlas through order seven and the
McKay catalogues for orders eight and nine \cite{mckay}. The verifier checks
coverage against the catalogues. Two implementations evaluated every
archived coefficient. A separate calculation using a different prime also
found nonzero coefficients in all $12\,111$ cases through order eight.
\end{proof}

\begin{center}
\begin{tabular}{lrrrrrrr}
\hline
Graph order & 3 & 4 & 5 & 6 & 7 & 8 & 9 \\
Connected classes verified & 2 & 6 & 21 & 112 & 853 & 11,117 & 261,080 \\
\hline
\end{tabular}
\end{center}

\newpage
\section*{Searching a graph beyond the certified range}
As a first order-ten check, the same method certifies all 317 connected
graphs with vertex degrees in $\{3,4\}$ and either 15 or 16 edges.

For each edge $e$, give its three entries variables $x_{e,0},x_{e,1},x_{e,2}$
with consecutive gaps at least one. Any bad lists can be sorted and rescaled
together to meet these constraints. For each index tuple $k\in\{0,1,2\}^E$,
impose the disjunction that some edge has equal endpoint sums under the
weighting $x_{e,k_e}$. These disjunctions exactly express that no allowed
weighting is proper. A search can add a disjunction whenever it finds a
proper weighting for the current candidate. Unsatisfiability of any subset
of disjunctions rules out bad lists on that graph.

For a real counterexample, choose one tied edge for each index tuple. The
tie equations have rational coefficients. Their simultaneous solution
space is defined over the rationals, and the strict-order constraints are
open within it. Thus a sufficiently close rational point remains a
counterexample. Clearing denominators gives integer lists and leaves
every witness tie intact. The unrestricted list conjecture remains open.

\section*{Reproducing the calculation}
The supplementary archive includes the input graph6 catalogues, the
coefficient table, and a checker. An exponent vector follows the
lexicographic ordering of the pairs of vertex labels in the input graph.
The checker reads the catalogue and the coefficient table in lockstep,
validates the exponent bounds, and recalculates the residue by direct
evaluation of the mixed difference. A zero residue from any choice of
exponents leaves the graph undecided and cannot certify a counterexample.

\begin{thebibliography}{9}
\bibitem{alon} N. Alon. Combinatorial Nullstellensatz. \emph{Combinatorics,
Probability and Computing} 8 (1999), 7--29.
\href{https://doi.org/10.1017/S0963548398003411}{doi:10.1017/S0963548398003411}.
\bibitem{bgn} T. Bartnicki, J. Grytczuk, S. Niwczyk. Weight choosability
of graphs. \emph{Journal of Graph Theory} 60 (2009), 242--256.
\href{https://doi.org/10.1002/jgt.20354}{doi:10.1002/jgt.20354}.
\bibitem{zhu} X. Zhu. Every nice graph is $(1,5)$-choosable.
\emph{Journal of Combinatorial Theory, Series B}, 2022.
\href{https://doi.org/10.1016/j.jctb.2022.08.006}{doi:10.1016/j.jctb.2022.08.006}.
\bibitem{mckay} B. D. McKay. \href{https://users.cecs.anu.edu.au/~bdm/data/graphs.html}
{Combinatorial data, simple connected graphs}.
\end{thebibliography}
\end{document}
